% GENERATED FILE. Edit canonical lessons, then run npm run generate:books.
% canonical-source-hash: fnv1a64:a29cef623f754150
% canonical-lessons: GU-C233-prani, GU-C233-ghodo, GU-C233-ut, GU-C233-balad, GU-C233-khedvu

\chapter{Animals that work}
\label{ch:gu-animals-that-work}
\hlchaptermodality{233}

\begin{chapteropening}
\textbf{By the end of this chapter:} \emph{I can say what the working animals did.}

Complete the last lesson of chapter 233: I can say what the working animals did.
\end{chapteropening}

\section[\texorpdfstring{\gu{પ્રાણી} (prāṇī)}{પ્રાણી (prāṇī)}]{\gu{પ્રાણી} (prāṇī) — animal, living creature}

\label{lesson:GU-C233-prani}



\begin{quote}
Before the new one: say the Gujarati for to take care, then the Gujarati for to get well.
\end{quote}

\subsection*{You'll want to know: \gu{પ્રાણી}}
\textbf{\gu{પ્રાણી}} — \emph{prāṇī} — \textquotedblleft{}animal, living creature\textquotedblright{}.

\begin{quote}
\textbf{\gu{મેં} \gu{જંગલમાં} \gu{ઘણાં} \gu{પ્રાણી} \gu{જોયાં}.} — \emph{mẽ jangalmā̃ ghaṇā̃ prāṇī joyā̃} — \textquotedblleft{}I saw many animals in the forest.\textquotedblright{}
\end{quote}

A neuter noun.

\begin{cousinweb}
\textbf{\gu{પ્રાણી}} (\emph{prāṇī}) is Sanskrit \emph{prāṇin}, \textquotedblleft{}living, a living creature\textquotedblright{}, built on \emph{prāṇa}, \textquotedblleft{}breath of life\textquotedblright{}, and taken over whole as a learned loan, a tatsama (CDIAL 8932).
\end{cousinweb}

\begin{grammarlens}[title={What you've built}]
One more word for everyday life at A2.
\end{grammarlens}

\subsection*{Your turn}
\begin{itemize}
\raggedright
  \item \emph{Say it:} \emph{prāṇī}
  \item \emph{Say it:} \emph{prāṇī}, once more
  \item \emph{Say it:} \emph{mẽ jangalmā̃ ghaṇā̃ prāṇī joyā̃} — I saw many animals in the forest.
  \item \emph{Recall:} say the Gujarati for to take care, then the Gujarati for to get well, then say \emph{prāṇī} again
\end{itemize}

\begin{tcolorbox}[breakable,colback=teal!4,colframe=teal!35!black,title={Before you move on}]
What does \gu{પ્રાણી} mean? (\textquotedblleft{}Animal, living creature\textquotedblright{}.) Say it once more.
\end{tcolorbox}
\section[\texorpdfstring{\gu{ઘોડો} (ghoḍo)}{ઘોડો (ghoḍo)}]{\gu{ઘોડો} (ghoḍo) — horse}

\label{lesson:GU-C233-ghodo}



\begin{quote}
Before the new one: say the Gujarati for to get well, then the Gujarati for animal.
\end{quote}

\subsection*{You'll want to know: \gu{ઘોડો}}
\textbf{\gu{ઘોડો}} — \emph{ghoḍo} — \textquotedblleft{}horse\textquotedblright{}.

\begin{quote}
\textbf{\gu{ઘોડો} \gu{સાંજે} \gu{ઘરે} \gu{આવ્યો}.} — \emph{ghoḍo sā̃je ghare āvyo} — \textquotedblleft{}The horse came home in the evening.\textquotedblright{}
\end{quote}

\begin{cousinweb}
\textbf{\gu{ઘોડો}} (\emph{ghoḍo}) is inherited from Sanskrit \emph{ghoṭa}, \textquotedblleft{}horse\textquotedblright{}, a word CDIAL takes to be of non-Aryan, probably Dravidian, origin (CDIAL 4516).
\end{cousinweb}

\begin{grammarlens}[title={What you've built}]
One more word for everyday life at A2.
\end{grammarlens}

\subsection*{Your turn}
\begin{itemize}
\raggedright
  \item \emph{Say it:} \emph{ghoḍo}
  \item \emph{Say it:} \emph{ghoḍo}, once more
  \item \emph{Say it:} \emph{ghoḍo sā̃je ghare āvyo} — The horse came home in the evening.
  \item \emph{Recall:} say the Gujarati for to get well, then the Gujarati for animal, then say \emph{ghoḍo} again
\end{itemize}

\begin{tcolorbox}[breakable,colback=teal!4,colframe=teal!35!black,title={Before you move on}]
What does \gu{ઘોડો} mean? (\textquotedblleft{}Horse\textquotedblright{}.) Say it once more.
\end{tcolorbox}
\section[\texorpdfstring{\gu{ઊંટ} (ū̃ṭ)}{ઊંટ (ū̃ṭ)}]{\gu{ઊંટ} (ū̃ṭ) — camel}

\label{lesson:GU-C233-ut}



\begin{quote}
Before the new one: say the Gujarati for animal, then the Gujarati for horse.
\end{quote}

\subsection*{You'll want to know: \gu{ઊંટ}}
\textbf{\gu{ઊંટ}} — \emph{ū̃ṭ} — \textquotedblleft{}camel\textquotedblright{}.

\begin{quote}
\textbf{\gu{ઊંટ} \gu{રણમાં} \gu{ચાલ્યો}.} — \emph{ū̃ṭ raṇmā̃ cālyo} — \textquotedblleft{}The camel walked in the desert.\textquotedblright{}
\end{quote}

\begin{cousinweb}
\textbf{\gu{ઊંટ}} (\emph{ū̃ṭ}) is a tadbhava. It comes from Sanskrit \emph{uṣṭra}, which meant a buffalo in the oldest texts and a camel later (CDIAL 2387).
\end{cousinweb}

\begin{grammarlens}[title={What you've built}]
One more word for everyday life at A2.
\end{grammarlens}

\subsection*{Your turn}
\begin{itemize}
\raggedright
  \item \emph{Say it:} \emph{ū̃ṭ}
  \item \emph{Say it:} \emph{ū̃ṭ}, once more
  \item \emph{Say it:} \emph{ū̃ṭ raṇmā̃ cālyo} — The camel walked in the desert.
  \item \emph{Recall:} say the Gujarati for animal, then the Gujarati for horse, then say \emph{ū̃ṭ} again
\end{itemize}

\begin{tcolorbox}[breakable,colback=teal!4,colframe=teal!35!black,title={Before you move on}]
What does \gu{ઊંટ} mean? (\textquotedblleft{}Camel\textquotedblright{}.) Say it once more.
\end{tcolorbox}
\section[\texorpdfstring{\gu{બળદ} (baḷad)}{બળદ (baḷad)}]{\gu{બળદ} (baḷad) — ox, bullock}

\label{lesson:GU-C233-balad}



\begin{quote}
Before the new one: say the Gujarati for horse, then the Gujarati for camel.
\end{quote}

\subsection*{You'll want to know: \gu{બળદ}}
\textbf{\gu{બળદ}} — \emph{baḷad} — \textquotedblleft{}ox, bullock\textquotedblright{}.

\begin{quote}
\textbf{\gu{બળદ} \gu{ખેતરમાં} \gu{ગયો}.} — \emph{baḷad khetarmā̃ gayo} — \textquotedblleft{}The ox went to the field.\textquotedblright{}
\end{quote}

\begin{cousinweb}
\textbf{\gu{બળદ}} (\emph{baḷad}) is a tadbhava, a word inherited from Sanskrit and reshaped by sound change. It comes from \emph{balivarda}, \textquotedblleft{}ox, bull\textquotedblright{} (CDIAL 9176).
\end{cousinweb}

\begin{grammarlens}[title={What you've built}]
One more word for everyday life at A2.
\end{grammarlens}

\subsection*{Your turn}
\begin{itemize}
\raggedright
  \item \emph{Say it:} \emph{baḷad}
  \item \emph{Say it:} \emph{baḷad}, once more
  \item \emph{Say it:} \emph{baḷad khetarmā̃ gayo} — The ox went to the field.
  \item \emph{Recall:} say the Gujarati for horse, then the Gujarati for camel, then say \emph{baḷad} again
\end{itemize}

\begin{tcolorbox}[breakable,colback=teal!4,colframe=teal!35!black,title={Before you move on}]
What does \gu{બળદ} mean? (\textquotedblleft{}Ox, bullock\textquotedblright{}.) Say it once more.
\end{tcolorbox}
\section[\texorpdfstring{\gu{ખેડવું} (kheḍvũ)}{ખેડવું (kheḍvũ)}]{\gu{ખેડવું} (kheḍvũ) — to plough, to till}

\label{lesson:GU-C233-khedvu}



\begin{quote}
Before the new one: say the Gujarati for camel, then the Gujarati for ox.
\end{quote}

\subsection*{You'll want to know: \gu{ખેડવું}}
\textbf{\gu{ખેડવું}} — \emph{kheḍvũ} — \textquotedblleft{}to plough, to till\textquotedblright{}.

\begin{quote}
\textbf{\gu{ગઈકાલે} \gu{તેણે} \gu{ખેતર} \gu{ખેડ્યું}.} — \emph{gaīkāle teṇe khetar kheḍyũ} — \textquotedblleft{}Yesterday he ploughed the field.\textquotedblright{}
\end{quote}

\begin{cousinweb}
\textbf{\gu{ખેડવું}} (\emph{kheḍvũ}) comes from \emph{*khēḍ-}, \textquotedblleft{}open up, plough\textquotedblright{}, a reconstructed form that CDIAL places under \emph{*khiḍ}, \textquotedblleft{}be open, expand\textquotedblright{}; it glosses the Gujarati verb \textquotedblleft{}to plough, drive\textquotedblright{} (CDIAL 3882).
\end{cousinweb}

\begin{grammarlens}[title={What you've built}]
One more word for everyday life at A2.
\end{grammarlens}

\subsection*{Your turn}
\begin{itemize}
\raggedright
  \item \emph{Say it:} \emph{kheḍvũ}
  \item \emph{Say it:} \emph{kheḍvũ}, once more
  \item \emph{Say it:} \emph{gaīkāle teṇe khetar kheḍyũ} — Yesterday he ploughed the field.
  \item \emph{Recall:} say the Gujarati for camel, then the Gujarati for ox, then say \emph{kheḍvũ} again
\end{itemize}

\begin{tcolorbox}[breakable,colback=teal!4,colframe=teal!35!black,title={Before you move on}]
What does \gu{ખેડવું} mean? (\textquotedblleft{}To plough, to till\textquotedblright{}.) Say it once more.
\end{tcolorbox}
